% !TEX program = xelatex
% =====================================================================
%  2026 校赛 B 题 智能球平衡控制装置 — 设计报告
%  编译：xelatex 设计报告.tex （连编两次以正确生成目录/交叉引用）
%  转 Word：pandoc 设计报告.tex -o 设计报告.docx
%  详见同目录 README.md / build.ps1
% =====================================================================
\documentclass[12pt,a4paper]{ctexart}

% ---------- 页面与版式 ----------
\usepackage[a4paper,top=2.5cm,bottom=2.5cm,left=3.0cm,right=3.0cm]{geometry}
\usepackage{setspace}
\onehalfspacing                      % 1.5 倍行距（电赛报告惯例）
\usepackage{indentfirst}             % 段首缩进

% ---------- 数学 / 符号 ----------
\usepackage{amsmath,amssymb}
\usepackage{bm}

% ---------- 图表 ----------
\usepackage{graphicx}
\usepackage{booktabs}                % 三线表
\usepackage{array}
\usepackage{multirow}
\usepackage{longtable}
\usepackage{float}
\usepackage{caption}
\captionsetup{font=small,labelsep=quad}
\captionsetup[figure]{name=图}
\captionsetup[table]{name=表}

% ---------- 列表 / 代码 ----------
\usepackage{enumitem}
\usepackage{xcolor}
\usepackage{listings}
\lstset{
  basicstyle=\ttfamily\small,
  breaklines=true,
  frame=single,
  rulecolor=\color{gray!50},
  backgroundcolor=\color{gray!5},
  numbers=left,
  numberstyle=\tiny\color{gray},
  showstringspaces=false,
  columns=fullflexible,
  keepspaces=true,
}

% ---------- TikZ 框图 ----------
\usepackage{tikz}
\usetikzlibrary{shapes.geometric,arrows.meta,positioning,calc,fit}
\tikzset{
  blk/.style={draw,rounded corners=2pt,align=center,minimum height=1.0cm,minimum width=2.6cm,font=\small},
  io/.style={draw,align=center,minimum height=0.9cm,minimum width=2.4cm,font=\small},
  dec/.style={draw,diamond,aspect=2,align=center,font=\footnotesize,inner sep=1pt},
  term/.style={draw,rounded corners=8pt,align=center,minimum height=0.8cm,minimum width=2.2cm,font=\small,fill=gray!8},
  sum/.style={draw,circle,minimum size=0.7cm,inner sep=0pt,font=\small},
  st/.style={draw,rounded corners=5pt,align=center,minimum height=0.9cm,minimum width=2.2cm,font=\small,fill=green!6},
  act/.style={draw,align=center,minimum height=0.8cm,minimum width=2.6cm,font=\footnotesize,fill=blue!4},
  ar/.style={-{Stealth[length=2mm]},thick},
}

% ---------- 超链接（放最后）----------
\usepackage[hidelinks]{hyperref}

% ---------- 章节编号对齐校赛格式：一、 / 1、 / （1） ----------
\ctexset{
  section/number=\chinese{section},
  section/name={,、},
  section/format=\Large\bfseries\raggedright,
  subsection/number=\arabic{subsection},
  subsection/name={,、},
  subsection/format=\large\bfseries,
  subsubsection/number={（\arabic{subsubsection}）},
  subsubsection/name={,},
  subsubsection/format=\normalsize\bfseries,
}

% ---------- 占位图框（照片未到时也能编译，后期替换 \includegraphics）----------
\newcommand{\figplaceholder}[2][12cm]{%
  \begin{center}
  \fbox{\begin{minipage}[c][4.5cm][c]{#1}\centering\color{gray} #2\end{minipage}}
  \end{center}}

% =====================================================================
\begin{document}

% --------------------------- 封面 ---------------------------
\begin{titlepage}
  \centering
  \vspace*{2.0cm}
  {\zihao{2}\bfseries 2026 年全国大学生电子设计竞赛（校赛）\par}
  \vspace{0.6cm}
  {\zihao{2}\bfseries 设\quad 计\quad 报\quad 告\par}
  \vspace{2.0cm}
  {\zihao{1}\bfseries 智能球平衡控制装置\par}
  \vspace{0.4cm}
  {\zihao{3}（B 题）\par}
  \vspace{2.5cm}
  {\zihao{4}
  \begin{tabular}{rl}
    题\quad 号： & B 题 \\[6pt]
    组\quad 别： & 高职组 \\[6pt]
    队\quad 员： & 张衡宇 \\[6pt]
    学\quad 校： & 深圳职业技术大学 \\[6pt]
    日\quad 期： & 2026 年 6 月 12 日 \\
  \end{tabular}\par}
  \vfill
\end{titlepage}

% --------------------------- 摘要 ---------------------------
\begin{center}{\zihao{3}\bfseries 摘\quad 要}\end{center}

本设计以 STM32H750 为主控，构建气流悬浮式乒乓球定高控制装置。系统采用超声波模块非接触测量球体高度，控制器以\textbf{悬停前馈 $u_{hover}$ + PD（球速阻尼）+ 弱积分}的控制律自动调节四线风机转速，使标准 40\,mm 乒乓球在透明导管内稳定悬浮于设定高度。针对竖直直筒"升力随高度衰减"构成的开环不稳定对象，引入球速微分反馈提供阻尼；针对风机低占空比段的强非线性，将 PWM 频率提升至 25\,kHz 以线性化执行特性、提高调节分辨率约 10 倍。人机交互采用 GC9A01 圆形彩屏，实时显示当前高度、目标高度、PWM 输出与控制状态，并以双色示波器曲线（目标参考线 + 实际高度）直观呈现定高与动态跟踪过程。系统上电就绪时间 $\le 2$\,s，支持 10/15/20\,cm 高度按键设定与定高控制。实测关键指标见正文第四节（部分指标待整定回填）。

\vspace{0.4cm}
\noindent\textbf{关键词：}\,悬浮定高；前馈+PD 控制；超声波测距；执行器线性化；STM32H750；GC9A01

\newpage
\tableofcontents
\newpage

% =====================================================================
\section{系统方案}
% =====================================================================

本系统由主控模块、测高模块、风机执行模块、人机交互模块与电源模块组成。被控对象为竖直透明导管内的标准 40\,mm 乒乓球，风机自底部送风使球悬浮，控制器实时检测球高并闭环调节风机转速以实现定高。下面分别论证各关键模块的方案选择。

\subsection{主控器件的论证与选择}

\textbf{方案一}：采用 51 系列单片机。成本低、开发简单，但主频与外设资源有限，难以同时承担高速定时器捕获测距、PWM 调速、彩屏刷新与实时控制运算。

\textbf{方案二}：采用 STM32F1 系列。资源较 51 充裕，可满足基本闭环，但浮点运算与高刷新彩屏驱动余量偏紧。

\textbf{方案三}：采用 STM32H750（Cortex-M7，480\,MHz，带 FPU 与大容量 SRAM）。高主频与硬件浮点利于 PD 控制与滤波运算，定时器与 SPI 资源充裕，可同时驱动 GC9A01 彩屏的示波器曲线。

综合主频、浮点能力与外设资源，且该平台工程已完成搭建并通过编译/烧录/串口/测距/显示的真机验证，\textbf{选择方案三，以 STM32H750 作为主控芯片}。

\subsection{测高模块的论证与选择}

\textbf{方案一}：红外测距模块。成本低、响应快，但易受环境光与管壁反射干扰，线性度差。

\textbf{方案二}：超声波测距模块（HY-SRF05）。非接触、受环境光影响小、成本适中，量程覆盖题目要求的 3\,$\sim$\,40\,cm，配合滤波可达 $\pm2$\,cm 量级精度。

\textbf{方案三}：TOF 激光测距模块。精度与刷新率高，可达 $\pm1$\,cm，但成本较高、驱动复杂。

综合成本、量程与精度需求，\textbf{基本要求阶段选择方案二（超声波）}；为冲击进阶要求的 $\pm1$\,cm，软件层以统一的"读高度"接口抽象底层传感器，预留\textbf{方案三 TOF 的平滑升级}而不改动控制逻辑。

\subsection{控制策略的论证与选择}

\textbf{方案一}：两段式 Bang-Bang（低于目标全功率、高于目标减力）。实现简单，但本质是开关控制，对无机械阻尼的悬浮对象必然产生极限环往复（真机实测已证实球大幅上下振荡、无法收敛）。

\textbf{方案二}：纯 PID。比例项可拉回目标，但竖直直筒近似开环不稳定，缺乏有效阻尼时易过冲、振荡。

\textbf{方案三}：悬停前馈 + PD（球速阻尼）+ 弱积分。前馈 $u_{hover}$ 直接托住悬停工作点，PD 的微分项对球速提供等效阻尼镇定对象，弱积分消除稳态静差。

综合稳定性与稳态精度，\textbf{选择方案三}，控制律见第二节式~\eqref{eq:law}。

\subsection{执行与调速方案的论证与选择}

\textbf{方案一}：1\,kHz PWM 驱动四线风机。真机实测升力区间被压缩在约 10 个占空比计数内（近似开关特性），控制被量化、误差大。

\textbf{方案二}：25\,kHz PWM 驱动四线风机（风机规格频率）。占空比与转速近似线性，有效调节区间显著展宽，分辨率提升约 10 倍。

\textbf{选择方案二}，由 TIM8 输出 25\,kHz PWM（引脚 PI6），占空比寄存器满量程 9600。

\subsection{显示方案的论证与选择}

\textbf{方案一}：单色 OLED。功耗低、驱动简单，但无法用颜色区分"目标曲线"与"实际曲线"。

\textbf{方案二}：GC9A01 1.28\,寸圆形彩屏。可用红/绿双色分别绘制目标参考线与实际高度曲线，直观呈现进阶要求的高度曲线显示。

\textbf{选择方案二}，以 SPI1 寄存器级直驱 GC9A01。

\subsection{系统总体方案}

系统总体框图如图~\ref{fig:block} 所示。控制链路为：超声波回波 $\rightarrow$ 距离换算 $\rightarrow$ 中值与滑动平均滤波 $\rightarrow$ 几何标定得球高 $\rightarrow$ 前馈+PD 闭环 $\rightarrow$ PWM 限幅/斜率限幅 $\rightarrow$ 风机；同时高度送 GC9A01 显示并经 UART 输出遥测以便在线整定。

\begin{figure}[H]
\centering
\begin{tikzpicture}[node distance=0.9cm and 1.4cm]
  \node[io] (sensor) {测高层\\HY-SRF05 超声波};
  \node[blk,right=of sensor] (mcu) {主控核心\\STM32H750\\(480\,MHz)};
  \node[io,right=of mcu] (fan) {执行层\\0420 四线风机\\TIM8 25\,kHz PWM};
  \node[io,below=of mcu] (hmi) {人机界面\\GC9A01 圆屏 + 4 按键};
  \node[io,above=of mcu] (dbg) {调试/遥测\\UART 串口};
  \draw[ar] (sensor) -- (mcu);
  \draw[ar] (mcu) -- (fan);
  \draw[ar] (mcu) -- (hmi);
  \draw[ar] (mcu) -- (dbg);
  \draw[ar,dashed] (fan.south) to[out=-90,in=0] node[below,font=\scriptsize]{气流抬升球体} (sensor.south);
\end{tikzpicture}
\caption{系统总体框图}
\label{fig:block}
\end{figure}

% =====================================================================
\section{系统理论分析与计算}
% =====================================================================

\subsection{被控对象建模}

设球质量 $m$、高度 $h$，受向上气流升力 $F(u,h)$ 与重力 $mg$，其中 $u$ 为风机 PWM 占空比。竖直导管内，球升高则远离风口、升力下降，即 $\partial F/\partial h<0$；增大占空比则升力上升，$\partial F/\partial u>0$。在悬停点 $(u_0,h_0)$ 附近线性化，记 $F_u=\partial F/\partial u$、$F_h=\partial F/\partial h$，运动方程为

\begin{equation}
m\ddot{h}=F(u,h)-mg\approx F_u\,\Delta u + F_h\,\Delta h .
\label{eq:plant}
\end{equation}

由式~\eqref{eq:plant} 得开环传递函数 $\dfrac{\Delta h(s)}{\Delta u(s)}=\dfrac{F_u}{ms^2-F_h}$。由于 $F_h<0$，其极点为 $s=\pm\sqrt{-F_h/m}$，一极点位于右半平面，对象\textbf{开环不稳定}且无固有阻尼。这从理论上解释了纯比例或 Bang-Bang 控制必然振荡的根因：必须人为引入阻尼项。

\subsection{控制律设计}

引入对球速 $\dot{h}$ 的微分反馈构成悬停前馈 + PD + 弱积分控制律，误差定义 $e=h_{\mathrm{target}}-h$：

\begin{equation}
u=u_{hover}+K_p\,e-K_d\,\dot{h}+K_i\!\int e\,\mathrm{d}t .
\label{eq:law}
\end{equation}

将式~\eqref{eq:law}（取 $K_i=0$ 分析镇定）代入式~\eqref{eq:plant}，闭环特征方程为

\begin{equation}
m s^2 + K_d F_u\, s + (K_p F_u - F_h)=0 .
\label{eq:char}
\end{equation}

由式~\eqref{eq:char}：当 $K_p>F_h/F_u$（恒成立，因 $F_h<0$）保证刚度为正，当 $K_d>0$ 引入正阻尼，全部极点即被拉入左半平面，系统镇定。微分项 $-K_d\dot{h}$ 是稳定的关键——这与真机现象一致：去掉 D 项球必发散冲顶或落底。等效阻尼比与无阻尼自然频率为

\begin{equation}
\omega_n=\sqrt{\frac{K_pF_u-F_h}{m}},\qquad
\zeta=\frac{K_dF_u}{2\sqrt{m\,(K_pF_u-F_h)}} .
\label{eq:zeta}
\end{equation}

整定目标取 $\zeta\approx0.6\sim0.8$（轻微过冲、快速稳定）。由于 $F_u,F_h$ 随工作高度变化、难以精确测量，工程上以式~\eqref{eq:zeta} 指导方向，最终用串口在线整定 $K_p/K_i/K_d$ 收敛到达标。闭环控制系统框图如图~\ref{fig:loop} 所示。

\begin{figure}[H]
\centering
\resizebox{0.95\linewidth}{!}{%
\begin{tikzpicture}[node distance=6mm and 10mm]
  \node[io] (ref) {目标高度\\$h_{target}$};
  \node[sum,right=of ref] (sum) {$\Sigma$};
  \node[blk,right=of sum] (ctl) {控制器\\$u_{hover}{+}K_pe{-}K_d\dot h{+}K_i\!\int\! e$};
  \node[blk,right=of ctl] (act) {执行器\\TIM8\\25\,kHz PWM};
  \node[blk,right=of act] (plant) {被控对象\\风机--导管--球};
  \node[right=8mm of plant] (out) {$h$};
  \node[blk,below=10mm of ctl,text width=3.0cm] (sen) {测高反馈\\超声波 + 3点中值 + 滑动平均};
  \draw[ar] (ref) -- node[above,font=\scriptsize]{$+$} (sum);
  \draw[ar] (sum) -- node[above,font=\scriptsize]{$e$} (ctl);
  \draw[ar] (ctl) -- node[above,font=\scriptsize]{$u$} (act);
  \draw[ar] (act) -- (plant);
  \draw[ar] (plant) -- (out);
  \draw[ar] (out.south) |- (sen.east);
  \draw[ar] (sen.west) -| node[pos=0.95,left,font=\scriptsize]{$-$} (sum.south);
\end{tikzpicture}}
\caption{闭环控制系统框图}
\label{fig:loop}
\end{figure}

\subsection{离散实现与抗饱和}

控制环以"测距新样本到达"为事件节拍驱动，步长 $\mathrm{d}t$ 取相邻样本真实间隔（标称 80\,ms，$12.5$\,Hz，受超声波最小触发间隔约 60\,ms 限制），避免固定节拍与传感器错相产生的锯齿速度。离散控制律如下：

\begin{lstlisting}[language=C]
e   = h_target - h;                       // 高度误差
v   = alpha*v + (1-alpha)*(h - h_prev)/dt;// 球速，EMA 平滑(alpha=0.6)
I   = clamp(I + e*dt, -I_lim, +I_lim);    // 条件积分抗饱和
u   = u_hover + Kp*e - Kd*v + Ki*I;       // 前馈 + PD + 弱I
u   = slew_limit(u, PWM_SLEW_PER_TICK);   // 斜率限幅防过冲
u   = clamp(u, 0, PWM_FULL_SCALE);        // 输出限幅
\end{lstlisting}

起调参数（集中于 \texttt{config.h}，待真机整定回填）：$K_p=150$、$K_i=8$、$K_d=120$，积分限幅 $\pm200$，误差死区 $\pm0.3$\,cm，斜率限幅 800\,/周期。

\subsection{执行器线性化的定量依据}

四线风机规格 PWM 频率为 25\,kHz。在 1\,kHz 时，占空比到转速严重非线性，真机实测整管升力被压缩在约 10 个占空比计数内（球在该窄区间内非顶即底，近似开关）。将定时器重配为 25\,kHz（预分频 0、自动重装 9599）后，占空比寄存器满量程由约 1000 提升至 9600，有效调节区间显著展宽，控制分辨率提升约 \textbf{10 倍}，前馈基准 $u_{hover}\approx3800$（约 $40\%$，待实测回填）。

\subsection{测高链路与几何标定}

声速法距离换算：$d=\dfrac{t_{\mathrm{echo}}}{58}$（cm，$t$ 单位 $\mu$s）。球高由几何关系（式~\eqref{eq:geo}）求得：

\begin{equation}
h=\mathrm{TUBE\_HEIGHT}-d_{\mathrm{raw}}-\mathrm{SENSOR\_OFFSET}-\delta_{\mathrm{ball}} ,
\label{eq:geo}
\end{equation}

其中腔体总高 46\,cm、传感器偏移 2\,cm、球径修正 4\,cm（均标 \texttt{[标定]}，需用钢尺在 10/20/30\,cm 三点校准回填）。滤波链：原始测距 3 点中值剔除尖刺 $\rightarrow$ 滑动平均（窗口 3，取小窗以降低相位滞后、保住 D 项质量）$\rightarrow$ 相邻跳变超 10\,cm 丢弃（连续丢弃 3 次强制接受防卡死）。

\subsection{误差预算}

测高误差预算见表~\ref{tab:err}。

\begin{table}[H]
\centering
\caption{测高误差预算}
\label{tab:err}
\begin{tabular}{lcl}
\toprule
误差源 & 估算 & 备注 \\
\midrule
超声波量化/分辨率 & $\pm0.3$\,cm & 58\,$\mu$s/cm，定时器计数足够 \\
管壁反射 / 气流扰动 & 【待实测】 & 中值+滑动平均后评估 \\
几何标定残差 & 【待实测】 & 三点标定后回填 \\
滤波相位滞后 & $\approx1$ 样本(80\,ms) & 窗口 3 折中 \\
\midrule
合计（目标） & $\le\pm1$\,cm & 题目基本 $\le\pm1$\,cm \\
\bottomrule
\end{tabular}
\end{table}

\subsection{时间预算}

系统就绪与实时性预算见表~\ref{tab:time}，就绪时间约 2\,s，满足题目"上电至正常工作 $\le5$\,s"。

\begin{table}[H]
\centering
\caption{时间预算}
\label{tab:time}
\begin{tabular}{lcl}
\toprule
阶段 & 耗时 & 说明 \\
\midrule
上电初始化 & 2000\,ms & 含 GC9A01 复位与初始化 \\
单控制周期 & 80\,ms & 测距$\rightarrow$滤波$\rightarrow$控制$\rightarrow$PWM \\
单帧曲线刷新 & $\approx60$\,ms & 184$\times$152 区域 @7.5\,MHz SPI \\
\midrule
就绪时间 & $\le2$\,s & $<$ 题目 5\,s 要求 \\
\bottomrule
\end{tabular}
\end{table}

% =====================================================================
\section{电路与程序设计}
% =====================================================================

\subsection{电路设计}

\subsubsection{主控与电源}
以 STM32H750XBH6 最小系统为核心，12\,V 电池供电，经降压为 MCU 提供 3.3\,V，风机直接由 12\,V 驱动。系统原理图见图~\ref{fig:sch}。

\begin{figure}[H]
\figplaceholder{此处插入系统原理图（导出后 \texttt{\textbackslash includegraphics} 替换占位框）}
\caption{系统电路原理图}
\label{fig:sch}
\end{figure}

\subsubsection{风机驱动（执行器）}
由 TIM8 输出 25\,kHz PWM（引脚 PI6）驱动四线风机的 PWM 控制线。四线风机引脚为 GND / +12\,V / PWM / TACH。其中 \textbf{TACH 转速反馈线经 10\,k$\Omega$ 上拉至 3.3\,V 接入 PC6（TIM3\_CH1，AF2）}，由定时器输入捕获测得脉冲周期换算转速（每转 2 脉冲，满速实测约 9520\,r/min），并随串口心跳输出（字段 \texttt{R}）。该转速反馈为后续风速串级闭环（发挥部分）提供内环测量量。

\subsubsection{超声波测距}
HY-SRF05 的 TRIG 由 GPIO 触发、ECHO 经定时器输入捕获测量回波脉宽，触发间隔 $\ge60$\,ms。

\subsubsection{显示接口（GC9A01）}
采用硬件 SPI1 寄存器级直驱（本工程不含 HAL SPI 库）。引脚映射见表~\ref{tab:pin}。

\begin{table}[H]
\centering
\caption{GC9A01 显示屏引脚映射}
\label{tab:pin}
\begin{tabular}{lll}
\toprule
GC9A01 & MCU 引脚 & 功能 \\
\midrule
SCL & PB3  & SPI1\_SCK (AF5) \\
SDA & PB5  & SPI1\_MOSI (AF5) \\
DC  & PG13 & GPIO 输出 \\
CS  & PE6  & GPIO 输出 \\
RST & PG14 & GPIO 输出 \\
\bottomrule
\end{tabular}
\end{table}

\subsection{程序设计}

\subsubsection{开发平台}
基于 STM32 HAL 库，使用 Keil MDK（ARM Compiler 6 / armclang）编译，SWD 烧录，并通过 UART 串口心跳输出遥测，配合上位机串口桥脚本在线整定参数（命令 \texttt{kp/ki/kd/uh} 等），无需反复烧录即可调参。

\subsubsection{软件架构}
按驱动层（Drivers）、算法层（Algorithm）、应用层（Core）三层分离，关键参数集中于 \texttt{config.h}：
\begin{itemize}[leftmargin=2em,itemsep=2pt]
  \item \textbf{Drivers}：超声波测距状态机、GC9A01 寄存器级 SPI 驱动、按键扫描、UART、风机转速捕获（TIM3）。
  \item \textbf{Algorithm}：前馈+PD 控制律、3 点中值 + 滑动平均滤波、球速估计（真实 $\mathrm{d}t$ + EMA）、条件积分抗饱和。
  \item \textbf{Core}：主循环调度器、UI 四界面状态机（HOME/CONTROL/CALIB/CURVE）与定高控制算法。
\end{itemize}

\subsubsection{主循环调度}
系统为非阻塞单循环调度：每轮依次执行非阻塞测距、按键与界面分发、串口命令解析、控制门判定，并按各自节拍刷新显示（200\,ms）与串口心跳（80\,ms），总体流程如图~\ref{fig:mainloop} 所示。控制环采用\textbf{事件驱动}——仅当测距产生新样本（\texttt{height\_updated}）时才执行一次闭环，使每步控制都用最新数据、$\mathrm{d}t$ 与采样真实间隔一致，避免固定节拍与传感器节拍错相产生锯齿球速。

\begin{figure}[H]
\centering
\begin{tikzpicture}[node distance=6mm and 14mm]
  \node[term] (s) {上电 / 外设初始化};
  \node[blk,below=of s] (meas) {非阻塞测距\\Ultrasonic\_Measure};
  \node[blk,below=of meas] (key) {按键扫描 + 界面分发\\KEY\_Process / UI\_OnKey};
  \node[blk,below=of key] (uart) {串口命令解析\\Process\_UART\_Command};
  \node[blk,below=of uart] (gate) {控制门（三分支，见图~\ref{fig:gate}）};
  \node[dec,below=8mm of gate] (dd) {到 200\,ms?};
  \node[act,right=of dd] (disp) {刷新界面\\OLED\_Update};
  \node[dec,below=10mm of dd] (dh) {到 80\,ms?};
  \node[act,right=of dh] (hb) {串口心跳\\H/T/E/P/M/D/R};
  \coordinate (L) at ([xshift=-12mm]meas.west);
  \draw[ar] (s)--(meas); \draw[ar] (meas)--(key); \draw[ar] (key)--(uart);
  \draw[ar] (uart)--(gate); \draw[ar] (gate)--(dd);
  \draw[ar] (dd)-- node[above,font=\scriptsize]{是} (disp);
  \draw[ar] (dd)-- node[left,font=\scriptsize]{否} (dh);
  \draw[ar] (disp) |- (dh);
  \draw[ar] (dh)-- node[above,font=\scriptsize]{是} (hb);
  \draw[ar] (hb) |- ([yshift=-7mm]dh.south);
  \draw[ar] (dh.south) -- node[left,font=\scriptsize]{否} ([yshift=-7mm]dh.south) -| (L) |- (meas.west);
\end{tikzpicture}
\caption{主循环调度流程}
\label{fig:mainloop}
\end{figure}

\subsubsection{控制门}
控制门按"是否使能、是否手动、是否有新样本"三级判定决定本轮动作，如图~\ref{fig:gate}。停机态直接丢弃挂起样本；手动标定态直给 PWM（用于扫 $u_{min}$/$u_{hover}$）；自动态仅在有新样本时跑一次闭环。

\begin{figure}[H]
\centering
\begin{tikzpicture}[node distance=7mm and 12mm]
  \node[dec] (e) {control\\\_enabled?};
  \node[act,right=of e] (e0) {丢弃挂起样本\\height\_updated=0};
  \node[dec,below=of e] (m) {manual\\\_mode?};
  \node[act,right=of m] (m1) {直给 manual\_pwm\\（跳过闭环）};
  \node[dec,below=of m] (h) {height\\\_updated?};
  \node[act,right=of h] (h0) {本轮不控制};
  \node[blk,below=of h] (c) {Control\_Update()\\前馈+PD+弱积分};
  \draw[ar] (e)-- node[above,font=\scriptsize]{否} (e0);
  \draw[ar] (e)-- node[left,font=\scriptsize]{是} (m);
  \draw[ar] (m)-- node[above,font=\scriptsize]{是} (m1);
  \draw[ar] (m)-- node[left,font=\scriptsize]{否} (h);
  \draw[ar] (h)-- node[above,font=\scriptsize]{否} (h0);
  \draw[ar] (h)-- node[left,font=\scriptsize]{是} (c);
\end{tikzpicture}
\caption{控制门三级判定}
\label{fig:gate}
\end{figure}

\subsubsection{超声波非阻塞测量}
测距采用两级非阻塞状态机：主状态 \texttt{measure\_state} 负责触发与回波等待，回波沿由 TIM4\_CH1 输入捕获中断（\texttt{capture\_state}：上升沿$\rightarrow$下降沿$\rightarrow$完成）测得脉宽。一次测量的处理流程如图~\ref{fig:ultra} 所示，含量程门限、3 点中值与跳变剔除（连续丢弃达上限则强制接受防卡死）。

\begin{figure}[H]
\centering
\begin{tikzpicture}[node distance=6mm and 10mm]
  \node[term] (idle) {空闲 (state=0)};
  \node[dec,below=of idle] (per) {到触发周期\\80\,ms?};
  \node[blk,below=of per] (trig) {发触发脉冲\\state=1};
  \node[dec,below=of trig] (cap) {回波完成?\\(capture=2)};
  \node[dec,right=14mm of cap] (to) {等待超时\\10\,ms?};
  \node[blk,below=of cap] (raw) {取距 $d=t_{echo}/58$};
  \node[dec,below=of raw] (rng) {$2{<}d{<}44$\,cm?};
  \node[blk,below=of rng] (med) {3 点中值滤波};
  \node[dec,below=of med] (jmp) {跳变${>}10$cm 且\\丢弃${<}3$次?};
  \node[blk,below=of jmp] (calc) {算高度 + 滑动平均\\height\_updated=1};
  \draw[ar] (idle)--(per);
  \draw[ar] (per)-- node[right,font=\scriptsize]{是} (trig);
  \draw[ar] (trig)--(cap);
  \draw[ar] (cap)-- node[above,font=\scriptsize]{否} (to);
  \draw[ar] (to.north) to[bend right=20] node[right,font=\scriptsize]{超时复位} (cap.east);
  \draw[ar] (cap)-- node[left,font=\scriptsize]{是} (raw);
  \draw[ar] (raw)--(rng);
  \draw[ar] (rng)-- node[left,font=\scriptsize]{是} (med);
  \draw[ar] (med)--(jmp);
  \draw[ar] (jmp)-- node[left,font=\scriptsize]{否} (calc);
\end{tikzpicture}
\caption{超声波非阻塞测量流程}
\label{fig:ultra}
\end{figure}

\subsubsection{定高控制算法}
\texttt{Control\_Update()} 为统一的前馈+PD+弱积分律，全程闭环（不再使用早期两段式 Bang-Bang）；\texttt{boost\_mode} 仅作显示标签（按 $|e|$ 远/近显示"起飞/定高"），不改变控制律。单次控制的计算流程如图~\ref{fig:ctrl}，含真实 $\mathrm{d}t$、误差死区、球速 EMA、条件积分抗饱和、输出限幅与 PWM 斜率限幅。

\begin{figure}[H]
\centering
\begin{tikzpicture}[node distance=5.5mm and 12mm]
  \node[term] (in) {新样本到达};
  \node[blk,below=of in] (dt) {真实 $\mathrm{d}t$（限幅 0.02$\sim$0.25\,s）};
  \node[blk,below=of dt] (err) {误差 $e=h_{target}-h$ + 死区};
  \node[blk,below=of err] (v) {球速 $v=(h{-}h_{prev})/\mathrm{d}t$，EMA};
  \node[blk,below=of v] (pd) {$u_{pd}=u_{hover}+K_pe-K_dv$};
  \node[dec,below=8mm of pd] (sat) {加积分后\\未饱和?};
  \node[act,right=16mm of sat] (frz) {冻结积分};
  \node[blk,below=8mm of sat] (intg) {积分累加 + 限幅};
  \node[blk,below=of intg] (out) {$u=u_{pd}+K_iI$，输出限幅};
  \node[blk,below=of out] (slew) {PWM 斜率限幅};
  \node[term,below=of slew] (wr) {写 TIM8\_CH2};
  \draw[ar] (in)--(dt); \draw[ar] (dt)--(err); \draw[ar] (err)--(v); \draw[ar] (v)--(pd); \draw[ar] (pd)--(sat);
  \draw[ar] (sat)-- node[above,font=\scriptsize]{否} (frz);
  \draw[ar] (sat)-- node[left,font=\scriptsize]{是} (intg);
  \draw[ar] (frz) |- (out);
  \draw[ar] (intg)--(out); \draw[ar] (out)--(slew); \draw[ar] (slew)--(wr);
\end{tikzpicture}
\caption{定高控制算法流程}
\label{fig:ctrl}
\end{figure}

\subsubsection{人机界面与曲线显示}
人机界面为四界面状态机，统一键义（K1 上/加、K2 进入/切换、K3 返回、K4 下/减），切换关系如图~\ref{fig:ui} 所示。CONTROL 界面同时显示当前高度、目标高度、PWM 输出、控制状态（对应数据显示要求）；CURVE 界面以红色参考线绘制目标高度、绿色曲线绘制实际高度，叠加 $\pm1$\,cm 误差带与 cm 刻度，直观呈现定高与动态跟踪。为降低刷新对控制环的阻塞，CURVE 采用扫描式渲染，每帧仅重绘并刷新最新一列。

\begin{figure}[H]
\centering
\begin{tikzpicture}[node distance=10mm and 16mm]
  \node[st] (home) {HOME 主菜单};
  \node[st,below left=12mm and 6mm of home] (ctrl) {CONTROL 定高};
  \node[st,below=14mm of home] (calib) {CALIB 标定};
  \node[st,below right=12mm and 6mm of home] (curve) {CURVE 曲线};
  \draw[ar] (home) to[bend right=12] node[left,font=\scriptsize]{K2} (ctrl);
  \draw[ar] (ctrl) to[bend right=12] node[right,font=\scriptsize]{K3} (home);
  \draw[ar] (home) to node[left,font=\scriptsize]{K2} (calib);
  \draw[ar] (calib) to[bend left=30] node[right,font=\scriptsize]{K3} (home);
  \draw[ar] (home) to[bend left=12] node[right,font=\scriptsize]{K2} (curve);
  \draw[ar] (curve) to[bend left=12] node[left,font=\scriptsize]{K3} (home);
  \node[font=\scriptsize,below=1mm of ctrl,align=center] {K1/K4 目标$\pm1$\\K2 切 10/15/20};
  \node[font=\scriptsize,below=1mm of curve,align=center] {同 CONTROL\\+ 实时曲线};
  \node[font=\scriptsize,below=1mm of calib,align=center] {K1/K4 PWM$\pm$步长\\K2 切步长 20/100/500};
\end{tikzpicture}
\caption{UI 四界面状态机}
\label{fig:ui}
\end{figure}

% =====================================================================
\section{测试方案与测试结果}
% =====================================================================

\subsection{测试条件与仪器}
测试条件：导管竖直固定、风道无人为干预、电源 12\,V。测试仪器：钢尺/卷尺（测高基准）、数字示波器（PWM 波形与 SPI 时序）、数字万用表、上位机串口调试工具。

\subsection{测试方案}

\subsubsection{基本要求测试}
\begin{itemize}[leftmargin=2em,itemsep=2pt]
  \item \textbf{高度测量}：CALIB 界面将球置于尺子 5/10/15/20/30\,cm，对比读数并观察抖动；判据误差 $\le\pm1$\,cm。
  \item \textbf{高度设定}：CONTROL 界面按键切换 10/15/20\,cm，核对屏显目标值。
  \item \textbf{定高控制}：分别设 10/15/20\,cm，记录稳定后 30\,s 内高度波动；判据稳态误差 $\le\pm2$\,cm。
  \item \textbf{数据显示}：核对当前高度/目标/PWM/状态四项是否同屏显示。
\end{itemize}

\subsubsection{进阶与发挥测试}
\begin{itemize}[leftmargin=2em,itemsep=2pt]
  \item \textbf{高精度定高}：升 TOF + 弱积分，测任意高度稳态误差 $\le\pm1$\,cm。
  \item \textbf{动态跟踪}：10$\leftrightarrow$20\,cm 切换，计进入并保持 $\pm1$\,cm 误差带的调节时间，判据 $\le5$\,s 且无明显振荡。
  \item \textbf{曲线显示}：核对坐标刻度、目标参考线、误差带与双色曲线。
  \item \textbf{抗扰}：吹/吸管口、扇风、敲管、瞬时遮挡传感器，观察恢复时间。
\end{itemize}

\subsection{测试结果}

\noindent\textbf{当前进度}：平台编译/烧录/串口/测距/显示已真机验证；GC9A01 彩屏点亮；风机能吹动球；定高闭环参数整定进行中。以下数据表待真机回填。

\begin{table}[H]
\centering
\caption{高度测量精度实测}
\label{tab:meas}
\begin{tabular}{cccc}
\toprule
尺子基准 & 装置读数 & 误差 & 达标 \\
\midrule
5\,cm  & 【待实测】 & & \\
10\,cm & 【待实测】 & & \\
20\,cm & 【待实测】 & & \\
30\,cm & 【待实测】 & & \\
\bottomrule
\end{tabular}
\end{table}

\begin{table}[H]
\centering
\caption{定高稳态误差实测}
\label{tab:hold}
\begin{tabular}{ccccc}
\toprule
目标 & 稳态均值 & 波动范围 & 误差 & 达标($\le\pm2$\,cm) \\
\midrule
10\,cm & 【待实测】 & & & \\
15\,cm & 【待实测】 & & & \\
20\,cm & 【待实测】 & & & \\
\bottomrule
\end{tabular}
\end{table}

\begin{table}[H]
\centering
\caption{动态跟踪调节时间实测}
\label{tab:track}
\begin{tabular}{cccc}
\toprule
切换 & 调节时间 & 是否振荡 & 达标($\le5$\,s) \\
\midrule
10$\rightarrow$20\,cm & 【待实测】 & & \\
20$\rightarrow$10\,cm & 【待实测】 & & \\
\bottomrule
\end{tabular}
\end{table}

\noindent 实测波形与照片（PWM 25\,kHz 波形、悬浮全景、CURVE 定高曲线截图）待补：
\begin{figure}[H]
\figplaceholder[10cm]{此处插入实测照片/截图}
\caption{实测现场与波形（待补）}
\label{fig:photo}
\end{figure}

% =====================================================================
\section{结论}
% =====================================================================

本设计以 STM32H750 为主控，针对竖直直筒开环不稳定对象，采用悬停前馈 + PD（球速阻尼）+ 弱积分控制律实现气流悬浮定高，并以 25\,kHz PWM 线性化执行器、以 GC9A01 圆屏双色示波器曲线呈现控制过程。系统就绪时间 $\le2$\,s，结构与算法已具备完成基本要求与冲击进阶要求的条件。

\noindent\textbf{创新点：}
\begin{enumerate}[leftmargin=2em,itemsep=2pt]
  \item 定位"升力区间被量化压窄"根因并以 25\,kHz 改造，将控制分辨率提升约 10 倍；
  \item 事件驱动的前馈+PD 闭环，以测距样本为节拍消除错相锯齿速度，针对双积分对象用球速阻尼镇定；
  \item 圆屏双色示波器曲线叠加 $\pm1$\,cm 误差带，使达标情况可视化自证。
\end{enumerate}

\noindent\textbf{不足与展望：}
\begin{enumerate}[leftmargin=2em,itemsep=2pt]
  \item 控制率受超声波限制为 12.5\,Hz，升级 TOF 可提升至 50\,Hz 以增大带宽；
  \item 风机 TACH 转速线未接入，后续补风速串级闭环可进一步提升抗扰与一致性。
\end{enumerate}

% =====================================================================
\section{参考文献}
% =====================================================================
\begin{thebibliography}{9}
\bibitem{rule} 全国大学生电子设计竞赛组委会. 2026 年校赛 B 题 智能球平衡控制装置题目[Z]. 2026.
\bibitem{rm} STMicroelectronics. STM32H750 Reference Manual RM0433[Z]. 2023.
\bibitem{ds} STMicroelectronics. STM32H750xB Datasheet[Z]. 2023.
\bibitem{auto} 胡寿松. 自动控制原理(第 7 版)[M]. 北京: 科学出版社, 2019.
\bibitem{astrom} Astrom K J, Murray R M. Feedback Systems: An Introduction for Scientists and Engineers[M]. Princeton University Press, 2008.
\bibitem{gc} Galaxycore. GC9A01 Datasheet[Z].
\bibitem{srf} HY-SRF05 超声波测距模块数据手册[Z].
\end{thebibliography}

% =====================================================================
\appendix
\renewcommand{\thesection}{附录\Alph{section}}
\ctexset{section/number=\Alph{section},section/name={附录,}}
\section{源代码目录结构}
\begin{lstlisting}
01_代码/
  config.h           # PID/PWM/几何/节拍集中配置
  Core/main.c        # UI + 控制状态机
  Drivers/           # 超声波 / GC9A01 / 按键 / UART
  Algorithm/         # 前馈+PD / 滤波 / 球速估计
\end{lstlisting}

\section{关键参数表（config.h）}
\begin{longtable}{lll}
\toprule
参数 & 取值 & 说明 \\
\midrule
\endhead
PWM\_FULL\_SCALE & 9600 & 25\,kHz 满量程 \\
PWM\_BASE ($u_{hover}$) & 3800 & 悬停前馈基准 [标定] \\
BOOST\_PWM & 5800 & 起飞推力 [标定] \\
$K_p,K_i,K_d$ & 150 / 8 / 120 & PID 起调值 [整定] \\
积分限幅 & $\pm200$ & 抗饱和 \\
误差死区 & $\pm0.3$\,cm & \\
控制周期 & 80\,ms & 12.5\,Hz \\
\bottomrule
\end{longtable}

\end{document}
